\begin{proposition}[One actual stage for the final family]\label{ex:app-final-family}
In the nonempty-seed case of Proposition~\ref{ex:app-nonempty-seeds},
fix the final $k,K$ and the stabilized common alphabet. For every
sufficiently large $b$, the address construction is realized at one
actual theta stage by a finite set $Q$ of distinct good odd primes
split in $K$, avoiding the fixed support and the specified finite
exclusions. With nonzero activations $\lambda_q$ it gives exactly
the family in (3.9),
\[
 h_x=\prod_{q\in Q}(q^*)^{\lambda_q(x)},\qquad h_0=1,
 \qquad q^*=(-1)^{(q-1)/2}q.
\]
All fixed local types are retained, and for every $x\ne0$,
\[
 \operatorname{ord}_{s=1}L(E^{h_x},s)=1,\quad L_*(h_xk)\ne0,
 \qquad j_{\rm gen}(h_x)\le\omega(h_x)+C_o,\quad
 v_2L_*(h_xk)\le2\omega(h_x)+C_e.
\]
The constants precede $b$. The finite graph budget and then the
actual theta precision are chosen after $b$. The theta prime remains
a coefficient prime and exclusion, not a relaxed Selmer prime.
\end{proposition}
\begin{proof}\leavevmode
For each chosen $b$, this construction has a finite budget $W_b$.
Enlarge it to bound every graph product and squarefree output in its
finite auxiliary tests, as well as every operator set. Contractions
only decrease the number of vertices. Thus every final $h_x$ lies in
the tested range of (F.7); this enlargement changes no fixed constant.
Only now take one actual theta stage in the finite intersection
of the selected table sets through $W_b+1$, the valid-stage set
for the full budget $|U|\le W_b$ with $R=2$, and the required
precision and finite exclusion sets.  In particular (F.6) holds
through the threshold $T$ used for that budget.
Realize its prescribed forward residue bits by distinct primes using CRT and Dirichlet;
quadratic reciprocity supplies the reverse bits.
There is no uniform bound on $W_b$, no intersection over all
budgets, and no product-surjectivity assumption on the alphabets.
The resulting one family satisfies (3.9), with constants fixed
before $b$. The theta prime $\ell_0$ is a coefficient prime and
an exclusion, not an active relaxed Selmer prime or a moving
arithmetic prime $r_i$.
\paragraph{Finite intersection and the upper bound.}
The required stage conditions form a finite list once the address
program is fixed. Let $\mathcal U$ be the chosen ultrafilter and
let $A_1,\ldots,A_t\in\mathcal U$ be the selected truth-table
sets used by that program, including all intermediate graphs.
Let $B\in\mathcal U$ impose (F.6) through the threshold $T$
for its operator budget and impose
$M\ge\max(0,\delta_o+W_b)+2$ and $M>T+\delta_o$.
Finite prime exclusions are included in $B$.
Then $B\cap\bigcap_{i=1}^t A_i\in\mathcal U$, hence is nonempty.
Choose a stage there before realizing the prime labels. The
coefficient rules hold for every realization at this stage, not
merely for a previously selected finite list of Fourier indices.
The address construction makes each nonzero address a unit of
both sources. For the odd source this says exactly
\[
 v_2a_M(h_x)=\delta_o+\omega(h_x)<M.
\]
The strict inequality follows from the tested-range precision
condition. Applying (O.8) gives the simple zero and
$j_{\rm gen}(h_x)\le\omega(h_x)+\delta_o+C_{\rm det}$.
The normalized even source gives
$v_2L_*(h_xk)\le2\omega(h_x)+C_{e,\rm final}$.
Thus one may take $C_o=\delta_o+C_{\rm det}$ and
$C_e=C_{e,\rm final}$, both fixed before $b$. Only the number
of tests, the threshold $T$, and the chosen stage depend on the
address dimension.
\end{proof}

